\pagebreak

\chapter{Accessing Specific FPGA Hardware Elements}\label{chap:asfpgahe}
    
    FPGA hardware elements can be created directly from 3PL source code.
    This can be done for the usual common elements that would normally be
    utilised during code generation or for little used or new elements that
    are not directly used by the code generator.
    
    {\bf Note that state elements, such as flip-flops and RAMs, created this way are
    outside the control of the normal execution thread mechanism.}
    
    There are
    three main reasons for using these directly created elements -
    \begin{itemize}
    \item   Special elements such as Virtex4 I/O serial/parallel converters
            cannot be directly incorporated into the semantics of the language.
    \item   We may wish to use low-level components whose inputs are
            not connected to a normal 3PL execution thread, for example a
            flip-flop whose clock enable or reset inputs are connected directly
            to some source.
    \item   We may wish to use absolute or relative placement of components,
            something which is very limited in 3PL proper.
    \item   The compiler may produce combinatorial logic which is sub-optimal
            and which in a critical case needs to be replaced by hand-coded
            logic.
    \end{itemize}
    In such cases we could wrap the created elements in a 3PL module,
    procedure or function to provide a normal 3PL code interface to hide
    these specialised constructs.
    
    As an example library io.3pl creates elements that are common to
    many FPGAs. These include IDDR, ODDR, IDELAY, IODELAY, IDELAYCTRL,
    ISERDES, SYSMON and some DSP functions. These are all created by the
    mechanism described below.
    
    Where an element is required which is not included in a library the
    following inbuilt procedures can be used to create it - \\
    {\bf element()} \autorefph{sec:ipelement}, \\
    {\bf elementend()} \autorefph{sec:ipelementend}, \\
    {\bf elementproperty()} \autorefph{sec:ipelementproperty}, \\
    {\bf elementinput()} \autorefph{sec:ipelmntin}, \\
    {\bf elementclockinput()} \autorefph{sec:ipelmntinclk} \\
    {\bf elementoutput()} \autorefph{sec:ipelmntout} \\
    {\bf elementtsoutput()} \autorefph{sec:ipelmntoutts}.
    
    The inbuilt procedure {\bf elementproperty()} only takes strings as
    value arguments so a number of library procedures are available
    in xilinx/common.3pl to format other types. These are - \\
    {\bf elementLogProperty()} \autorefph{sec:lpelementLogProperty}, \\
    {\bf elementIntProperty()} \autorefph{sec:lpelementIntProperty}, \\
    {\bf elementBinProperty()} \autorefph{sec:lpelementBinProperty}, \\
    {\bf elementHexProperty()} \autorefph{sec:lpelementHexProperty}, \\
    {\bf elementFloatProperty()} \autorefph{sec:lpelementFloatProperty}, \\
    {\bf elementStrProperty()} \autorefph{sec:lpelementStrProperty}
    
    As an example if we wished to directly use a CLB shift register instead of
    using using it implicitly with {\bf del()}, we could create it by -
    \begin{lstlisting}[gobble=4, language=threepl]
       string(id, "srl123");
       element("SRL16E", id);
    \end{lstlisting}
    where
    "SRL16E" is the name of the element in the Xilinx library and "srl123" is an
    optional identifier for the element.
    We would now connect up each of its ports -
    \begin{lstlisting}[gobble=4, language=threepl]
       value({d, ce, q}, "log");
       static(a, "uint:4");
       d := .....;
       a <- .....;

       elementinput(id, "D", d);
       elementinput(id, "A", a, 1);
       elementclockinput(id, "CLK", currentclock());
       elementinput(id, "CE", ce);
       elementoutput(q <- id, "Q");
       elementend(id);
       .... <- q;
       .
       .
    \end{lstlisting}
    We could have omitted the {\bf id} argument in all the calls above in which
    case the currently open element would have been correctly implied.

    
    The call to {\bf element()} creates the element which is retained until the
    next call to {\bf element()} as the current element. Calls to {\bf
    elementinput()}  or {\bf elementoutput()} with no {\bf id} apply to the
    current element. When an id is explicitly used then multiple elements can
    be created using interleaved calls to the above procedures, though there is
    probably little advantage in doing that. The arguments to {\bf
    elementinput()} are -
    \begin{itemize}
    \item   the optional element id (string)
    \item   the port name (string)
    \item   the signal variable
    \item   an array format indicator (integer 0 for single signal, 1 for an array
            where the pins are specified individually, e.g.
            {\bf (port A0 (direction INPUT))} or 2 for an EDIF
            array format such as {\bf (port (array (rename A "a[0:11]") 12) (direction INPUT)))}.
            Unfortunately the required format seems to depend on whether ISE or
            Vivado are used and also on the specific element.
    \end{itemize}
    Input signal variables can be the constants true or false, any value,
    selectvalue, static or priority variable or an expression. The signal cannot be
    from a queue read or an expression containing a queue read but can be from a queue
    variable if the examine operator \verb'?' is used.
    
    This process is subject to the following observations -
    \begin{enumerate}
    \item   The variables or constants are not checked against the port
            type for validity, e.g. if an output port was connected to a
            constant or a static (i.e. the output of the static) this
            would not be detected until code generation time when a fatal
            error resulting from multiple net sources would occur.
    \item   Errors may not be detected until the Xilinx place-and-route
            software is run at which point an obscure error message
            would be generated.
    \item   Elements generated this way cannot be handled by the simulator
            (howevere the simulator is obsolete and not currently available).
    \end{enumerate}

    In the example above it would have been sensible to construct a module or
    procedure so that the code could be reused, for example - 
    \begin{lstlisting}[gobble=4, language=threepl]
       module
       srl16 (value(out) <- value(in), value(ce, "log"), value(addr, "uint:4")) {
           element("SRL16E");
           elementinput(, "D", in);
           elementinput(, "A", addr, 1);
           elementclockinput(, "CLK", currentclock());
           elementinput(, "CE", ce);
           elementoutput(out <- , "Q");
           elementend();
       }
    \end{lstlisting}
    It makes no difference if this is a module or procedure as there is no in-line
    code and no queue reads or writes and the calling syntax and argument passing
    for these are identical.
    
    To implement a similar delay for say an 8-bit integer -
    \begin{lstlisting}[gobble=4, language=threepl]
       module
       srl16_8 (value(out, "int:8") <- value(in, "int:8"), value(ce, "log"), value(addr, "uint:4")) {
           value(o, "[8]bits:1");
           for (int(i) ; i<8 ; i++)
               srl16(o[i] <- getbit(in, i), ce, addr); 
           out := cast(o, "int:8");
       }
    \end{lstlisting}
    To access single bits of the input the library function getbit() (stdlib.3pl)
    has been used. There is no equivalent for assigning bits to the output so a
    local array of single bits has been defined and then this was cast to the
    output type.
    
    The following code is a simplified example of the use of such a delay. To avoid
    complicating the example the inputs are assumed to be synchronous with the
    current clock. 
    \begin{lstlisting}[gobble=4, language=threepl]
       iopins(io_pins, "[20]str", {"AB12", "AC13", .......});
       
       input(din, "int:8", loc=io_pins[0..7], readdomain=null);
       input(dvalid, "log", loc=io_pins[8], readdomain=null);
       static(sum, "int:8");
       value(del, "int:8");
       value(step, "log");

       srl16_8(del <- din, step, 5);

       when (dvalid) par {
           sum <- din + del;
           assert(step <-);
       }

       output(,, sum, loc=io_pins[9..16]);
    \end{lstlisting}
    
    In the above example the clock enable is sourced from an instance of the {\bf
    assert()} \autorefph{sec:ipassert} inbuilt procedure. This procedure asserts the
    output when the execution stream in which it is contained executes.
    
    Alternatively the SRL16E CE could have been driven directly from variable
    {\bf dvalid}.
    
    \begin{lstlisting}[gobble=4, language=threepl]
       .
       .
       srl16_8(del <- din, dvalid, 5);

       when (dvalid)
           sum <- din + del;
       .
       .
    \end{lstlisting}
